\chapter{Specifying events of interest}
\label{ch:specify}

The first stage of our approach addresses \emph{what} weather scenario should
be generated. We specify an event through two components: a mask $\pi$ that
selects the atmospheric variable and spatial region of interest, and a target
trajectory (TODO: insert y here with correct notation) describing the desired evolution of the corresponding regional
quantity over the guided forecast horizon.

Throughout this thesis, we focus on regional surface-temperature heatwaves.
Accordingly, $\pi$ selects surface temperature over a prescribed region, while
the target trajectory specifies how its regional average should evolve over
the guided weather steps. The formulation is nevertheless not restricted to
temperature and can be applied to other variables or (TODO: reformulate) differentiable regional
functionals of the forecast state.

(TODO: improve introductory sentence)
As generative model, ArchesWeather weather can generate ensemble forecasts by drawing multiple samples.
To generate weather scenarios we want to specify the guidance target 
to retain the properties of the ensemble. Rather than guiding all
members toward the same absolute trajectory, we define a common relative
displacement and apply it to each member with respect to its matched unguided
forecast. This preserves the member-specific structure of the target
trajectories instead of imposing an artificial collapse onto a single value.

We first define the target variable and region through the spatial mask
$\pi$. We then construct the target trajectory over the guided horizon and
derive the corresponding member-specific targets to use in the guidance
algorithm.

\section{Target variable and region}

We jointly specify the target variable and spatial region through a mask
$\pi$ over the weather state. Let

\[
    x_n \in \mathbb{R}^{V\times I\times J}
\]

denote a weather state, where $V$ indexes atmospheric fields, while $I$ and
$J$ index latitude and longitude. We denote the corresponding spatial
coordinates by

\[
    (\varphi_i,\psi_j),
    \qquad
    i=1,\ldots,I,
    \quad
    j=1,\ldots,J.
\]

Let $v^\star$ denote the target variable, and let

\[
    k_{ij}\in\{0,1\}
\]

indicate whether spatial grid cell $(i,j)$ lies inside the selected target
region. In this thesis, the spatial support is defined through a rectangular
bounding box on the latitude--longitude grid.

The mask is non-zero only for the selected variable and within the selected
region. To obtain an area-weighted regional average, each included grid cell
is weighted according to the physical surface area that it represents. For a
regular latitude--longitude grid, the area contribution of latitude row $i$ is
proportional to

\begin{equation}
    a_i
    =
    \int_{\beta_i^-}^{\beta_i^+}
    \cos\varphi\,\mathrm{d}\varphi,
    \label{eq:mask-cell-area}
\end{equation}

where $[\beta_i^-,\beta_i^+]$ denotes the latitude extent of the corresponding
grid cell. Since the longitudinal grid spacing is constant, its contribution
cancels during normalization.

We therefore define the normalized mask as

\begin{equation}
    \pi_{vij}
    =
    \frac{
        \mathbf{1}[v=v^\star]\,
        k_{ij}\,a_i
    }{
        \displaystyle
        \sum_{i',j'}
        k_{i'j'}\,a_{i'}
    },
    \label{eq:target-mask}
\end{equation}

such that

\[
    \sum_{v,i,j}\pi_{vij}=1.
\]

The corresponding masked-average operator is

\begin{equation}
    \mathcal{M}_\pi(x_n)
    :=
    \sum_{v,i,j}
    \pi_{vij}\,x_{n,vij}.
    \label{eq:masked-average}
\end{equation}

For the experiments considered in this thesis,
$\mathcal{M}_\pi(x_n)$ therefore represents the area-weighted mean surface
temperature within the selected target region.

Area weighting prevents the regional average from being biased toward
high-latitude grid cells, which represent a smaller physical area on the
sphere. This follows the latitude-weighting convention used in
WeatherBench~2 \parencite{rasp2024weatherbench2}.

Figure~\ref{fig:target-mask-example} illustrates the resulting area-weighted
mask for an example target region over South America. Its spatial support is
rectangular on the latitude--longitude grid, while the normalized weights vary
with latitude according to the physical area represented by each grid cell.

TODO: improve the figure by having a single colorbar on the side and the globe has the same diameter as the side of the rectangle.
\begin{figure}[t]
    \centering
    \begin{subfigure}[c]{0.48\textwidth}
        \centering
        \includegraphics[height=5.2cm]{figures/mask_flat.png}
        \caption{Latitude--longitude representation.}
        \label{fig:target-mask-flat}
    \end{subfigure}
    \hfill
    \begin{subfigure}[c]{0.48\textwidth}
        \centering
        \includegraphics[height=5.2cm]{figures/mask_globe.png}
        \caption{Spherical representation.}
        \label{fig:target-mask-globe}
    \end{subfigure}

    \caption{
        Example of the spatial target mask. The rectangular support determines
        which grid cells are included, while the latitude-dependent weights
        account for differences in physical grid-cell area.
    }
    \label{fig:target-mask-example}
\end{figure}

Using a binary spatial support makes the target region explicit. Within the
selected region, grid cells differ only through their physical area and are
otherwise treated uniformly. This avoids introducing an additional spatial
preference, as would occur, for example, with a bell-shaped mask whose center
is weighted more strongly than its boundaries. Normalizing $\pi$ to sum to one
ensures that $\mathcal{M}_\pi$ remains an average and is therefore directly
comparable across regions of different sizes.

Although $\pi$ selects one particular variable and region, its spatial support
$k_{ij}$ can be reused to construct analogous regional averages for other
atmospheric fields. This becomes useful when evaluating how guidance applied to
surface temperature affects other variables and pressure levels over the same
region.

\section{Target trajectory}
\label{sec:target-trajectory}

Once the target variable and region have been specified through the mask
$\pi$, the remaining task is to define the desired evolution of the targeted
quantity. Let $H\leq N$ denote the final guided weather step. We first describe
this evolution at the ensemble level through the target trajectory

\[
    y^\star_{1:H}
    =
    \left(y^\star_1,\ldots,y^\star_H\right),
\]

and subsequently translate it into member-specific targets
$y^\star_{1:H,m}$. No target is imposed for $n>H$.

TODO: make sure that $N$ is actually defined in the background, it also doesn't hurt to respecify it here. 

\subsection{Guiding an ensemble}

At each weather step, ensemble members generally attain different values of
the targeted quantity $\mathcal{M}_\pi(x_n)$, with their spread typically
increasing as forecast uncertainty grows with lead time.

Assigning the same absolute target trajectory to every member would ignore
these differences and prescribe zero ensemble spread in
$\mathcal{M}_\pi(x_n)$. (TODO: check wording) We therefore define the guidance target as desired change relative to each
member's own matched unguided trajectory.

Each guided rollout is paired with a matched unguided rollout

\[
    \left\{
        x_{n,m}^{\mathrm{ung}}
    \right\}_{n=1}^{N},
    \qquad
    m=1,\ldots,M,
\]

using the same initial atmospheric context and matched stochastic realizations.
At weather step $n$, we prescribe a common relative change $\rho_n$ across the
ensemble and define the target of member $m$ as

\begin{equation}
    y_{n,m}^\star
    :=
    (1+\rho_n)\,
    \mathcal{M}_\pi
    \left(
        x_{n,m}^{\mathrm{ung}}
    \right),
    \qquad
    n=1,\ldots,H.
    \label{eq:member-target}
\end{equation}

The corresponding ensemble-level target is the mean of the member-specific
targets,

\begin{equation}
    y_n^\star
    :=
    \frac{1}{M}
    \sum_{m=1}^{M}
    y_{n,m}^\star
    =
    (1+\rho_n)
    \frac{1}{M}
    \sum_{m=1}^{M}
    \mathcal{M}_\pi
    \left(
        x_{n,m}^{\mathrm{ung}}
    \right).
    \label{eq:ensemble-target}
\end{equation}

Applying the same relative change to all members retains the differences
between their matched unguided trajectories. Let
$\operatorname{std}_m[\cdot]$ denote the standard deviation across the
ensemble-member index $m=1,\ldots,M$ at a fixed weather step $n$. Then, for
$1+\rho_n>0$, the spread of the member-specific targets satisfies

\begin{equation}
    \operatorname{std}_m
    \left[
        y_{n,m}^\star
    \right]
    =
    (1+\rho_n)\,
    \operatorname{std}_m
    \left[
        \mathcal{M}_\pi
        \left(
            x_{n,m}^{\mathrm{ung}}
        \right)
    \right].
    \label{eq:target-ensemble-spread}
\end{equation}

The target ensemble therefore inherits the lead-time-dependent spread of the
matched unguided ensemble, scaled by the prescribed relative change, rather
than collapsing all members onto a common trajectory. Whether the guided
ensemble reflects the spread implied by these member-specific
targets ultimately depends on the dynamics of the guidance algorithm. 


\subsection{Prescribing the relative change}

It remains to specify the relative change $\rho_n$. In principle, the shape
and magnitude of $\rho_{1:H}$ could be informed by historical events,
climatological statistics, or other application-specific definitions of the
scenario of interest. Historical events could, for example, provide a
reference for how the requested relative intensity evolves over time.

(TODO: state this also in the intro for consistency)
The focus of this thesis, however, is to evaluate the behaviour of the guidance
method under controlled settings. We therefore prescribe $\rho_n$ directly and
use a simple linear increase over the guided horizon,

\begin{equation}
    \rho_n
    :=
    \frac{n}{H}\rho_H,
    \qquad
    n=1,\ldots,H,
    \label{eq:relative-target-schedule}
\end{equation}

where $\rho_H$ denotes the maximum requested relative change and $H$ the final
guided weather step.

This construction gradually increases the requested deviation rather than
imposing its full magnitude from the first forecast step. Because the relative
change is applied to each member's matched unguided value through
\Cref{eq:member-target}, the target trajectory continues to inherit the
individual temporal evolution of that member while progressively moving away
from its unguided forecast. At step $H$, the prescribed relative change reaches
$\rho_H$ and guidance is subsequently stopped.

For the surface-temperature experiments considered here, $\rho_H$ is therefore
a scale-dependent experimental intensity parameter rather than a climatological
measure of heatwave extremeness.
TODO: specify this how we will use it in our experiments (we will define multiple levels of intensity and assess results under the different levels)


\section{Pipeline overview}
In practical terms, the elements specified throughout this chapter are changed together as follows.
Starting from an initial atmospheric context
$(x_{-1},x_0)$, we first define the mask $\pi$, which specifies the target
variable and spatial region. For each ensemble member $m=1,\ldots,M$, we then
generate a matched unguided rollout

\[
    \left\{
        x_{n,m}^{\mathrm{ung}}
    \right\}_{n=1}^{N},
\]

which serves as the reference trajectory for constructing the target.

A relative intensity trajectory $\rho_{1:H}$ is prescribed over the guided
horizon and applied to each matched unguided member according to
\Cref{eq:member-target}. This yields the set of member-specific target
trajectories

\[
    \left\{
        y_{1:H,m}^\star
    \right\}_{m=1}^{M}.
\]

The output of the specification stage is therefore

\begin{equation}
    \left(
        \pi,\,
        \left\{
            y_{1:H,m}^\star
        \right\}_{m=1}^{M}
    \right),
    \label{eq:specification-output}
\end{equation}

where $\pi$ determines which atmospheric quantity is targeted, while the
member-specific trajectories determine its desired evolution during the guided
forecast horizon. For $n>H$, no target is imposed and the forecast continues
without guidance until the end of the rollout at $N$.

TODO: include here the rollout trajectories planning figure.

These objects define \emph{what} the model should generate independently of
\emph{how} the generative process is modified to realize them. The next chapter
addresses this complementary guidance problem.

% TODO: Add a figure illustrating the resulting ensemble of target trajectories,
% together with the corresponding matched unguided trajectories.

TODO: somewhere around here we should show the graph that this creates.

\includegraphics{figures/rollout-trajectories.png}